\chapter{生きた細胞の計算機}
% full version: chapters/21_living.tex

\pencilfig{21}{\pencilart{21}}{チップの上の生きたニューロン：その下には数千の電極があり、どれも聞くことも話すこともできる。}

シリコンではなく生きたニューロンでマシンを作れるだろうか。作れるし、すでに作られている。ニューロンは、数千の電極をもつチップの上で直接育てられる。平らな層として、あるいは小さな脳に似た塊として。どの電極も、聞くことも信号を送ることもできる。これは回路が丸見えのテストベンチであり、何でも測定でき、何にでも介入できる。ただし放送局はない。制御レジスタのないデータバスである。

\section*{ネットワークが自分ですること}

放っておくと、培養したネットワークはときどき全体が一斉に燃え上がる。ほぼすべてのニューロンが一緒にカチッと鳴り、それから静まる。その間隔は1秒から数分だ。これはおなじみの発振器である。ネットワークが自分で自分を興奮させ、接続部が物質の蓄えを回復するまで疲れているのだ。

\pencilfig{129}{%
% spike raster of 8 neurons in a culture left to itself: sparse single clicks and shared bursts
\foreach \b in {0.9,3.1,4.0,6.6}{\fill[pcAmber, opacity=0.18] (\b-0.1,-0.15) rectangle (\b+0.32,2.95);}
\foreach \y/\ss in {0/{0.96,1.0,1.13,2.43,3.0,3.12,3.24,4.02,4.09,4.15,6.66,6.69,6.81},
  1/{0.46,0.88,1.0,1.05,3.09,3.2,3.94,4.04,4.37,6.65,6.71},
  2/{0.73,0.84,0.94,1.41,3.08,3.12,3.21,4.05,4.06,4.17,6.59,6.63},
  3/{0.89,0.93,1.06,3.1,3.19,3.28,3.89,3.94,4.13,4.2,4.31,6.63,6.65,6.78},
  4/{0.87,0.96,3.09,3.15,3.23,3.73,3.96,4.06,4.19,5.98,6.66,6.68,6.75},
  5/{0.44,0.92,0.97,3.05,3.22,3.27,4.01,4.07,4.17,5.76,6.57,6.73,6.75},
  6/{0.92,1.0,1.73,2.85,3.18,3.19,4.0,4.07,6.53,6.66,6.73},
  7/{0.87,0.92,3.13,3.13,3.27,3.99,4.06,4.17,6.44,6.61,6.72,7.13}}{
  \foreach \s in \ss {\draw[pcBlue, line width=0.7pt] (\s,0.35*\y) -- (\s,0.35*\y+0.25);}}
\draw[arr] (0,-0.3) -- (7.5,-0.3); \node[lbl, anchor=north] at (3.75,-0.32) {時間};
\node[lbl, anchor=east, align=right] at (-0.05,1.35) {8個の\\ニューロン};
\node[lbl, text=pcAmber!70!black, anchor=south] at (3.55,2.95) {一斉の燃え上がり};
}{放っておかれたチップ上の培養。単発のカチッはまれで、ときどきネットワークのほぼ全体が一斉に燃え上がり、接続部が蓄えを回復するまで静まる。}

\section*{ドーパミンのない教師}

ドーパミンがないのに、こうしたネットワークをどう教えるのか。最初の方法は電気的なものだった。たとえば、ネットワークが望みの応答を出すまで刺激し、出したらすぐに止める。すると応答が定着する。別の実験では、培養が単純な仮想の「体」を操り、それを目標へ導くことを学んだ。これらの実験ではどれも、教師はエンジニアが選ぶ電流のパターンである。

\section*{光の閃光でドーパミンを}

あるとき、教師を化学的なものにする試みがなされた。生きたニューロンの塊を使うあるプラットフォームで、培養液に、ドーパミンの作用を封じる光感受性の「かご」に入れたドーパミンが加えられた。ネットワークが刺激に以前より強く応答すると「ご褒美」が与えられた。紫外線の閃光がかごを壊し、ドーパミンが外に出るのだ。240回の繰り返しで、応答は全体として大きくなった。だが著者ら自身が正直に書いているように、この1回の実験だけでは、原因がドーパミンだとは言えない。

\section*{台車の上の棒}

最近の実験では、ニューロンの塊に台車の上で棒のバランスをとることを教えた。学習するプログラムの定番課題である。どの教育用パルスを与えるかはプログラムが選んだ。プログラムが選んだ信号はランダムな信号よりよく効いたが、45分休ませると、学んだことは消えてしまった。そして\nt{GLUT}接続部を遮断すると、学習はまったく進まなかった。学んだことはまさにそこに宿っているのだ。一方、教師は依然として外にいる。ただ今度はエンジニアではなく、プログラムである。

\section*{生きたリザーバ}

別のアプローチもある。生きたネットワークをまったく教えないのだ。それを「リザーバ」として使う。豊かで入り組んだシステムであり、その応答を単純な電子式の読み取り器が解読することを学ぶ。こうして生きた塊が声の聞き分けを助けた。ただし著者らによれば、動作中に塊そのものの結合も変わっていた。

\section*{コストと問い}

チップ上の100万個のニューロンが消費するのは約1万分の1ワットで、周りの電子回路よりずっと少ない。そしてテストベンチは、何を成功とみなすかを自分で決められない。その役は常に外の誰かが担う。さらに、こうしたシステムが大きくなるにつれて、倫理の問題も浮かび上がる。それはまだ誰も解決していない。

\fullversion{21}
